\chapter{Final Capstone and Certification Preparation}
\index{certification}\index{Data Engineer Associate}\index{Data Engineer Professional}%
\index{capstone project!final}\index{mock exam}\index{architecture review}\index{study plan}%
\begin{objectives}
\begin{itemize}
    \item Integrate all six milestones into one production-grade, code-deployed platform.
    \item Execute the clean-environment redeploy drill and the replay drill as final evidence.
    \item Map the Databricks certification domains to this book's chapters and labs.
    \item Run two timed mock exams and remediate weak domains by rebuilding labs.
    \item Present the platform as a 15-minute architecture defense with numbers.
    \item Package the six capstones into a portfolio that survives an interviewer's probing.
\end{itemize}
\end{objectives}

\begin{crosscloud}
Certification week looks different per platform, but the \textbf{readiness method} is identical:
domains mapped to labs, mocks under time, remediation by rebuilding.

\vspace{2mm}
{\footnotesize
\begin{tabular}{@{}p{3.0cm}p{3.4cm}p{3.2cm}p{3.2cm}@{}}
\toprule
\textbf{Element} & \textbf{Databricks} & \textbf{AWS} & \textbf{Azure / GCP / Fabric} \\
\midrule
Target exam & Data Engineer Associate / Professional & Data Engineer Associate & DP-700 / Professional Data Engineer / DP-600 \\
Heavy domains & Platform, Delta Lake, Spark ETL, pipelines, governance & Storage, ingestion, transform, ops & Fabric workloads / GCP services / Power BI \\
Lab weight & High (build fluency) & High & High \\
Mock standard & Two mocks, 80\%+, no domain $<$70\% & Same & Same \\
Capstone defense & Architecture review with cost numbers & Same & Same \\
\bottomrule
\end{tabular}}

\vspace{1mm}
The portability argument of the whole series: if you can defend the Databricks capstone \emph{and}
sketch its AWS or Fabric equivalent from the cross-cloud boxes, you have learned data engineering,
not one console.
\end{crosscloud}

\section{Week 24 Roadmap and Mental Model}
Week 24 is integration and examination. The capstone reassembles the program---lakehouse, pipelines, streaming, ML, governance, delivery---into one platform deployed from one repository. The certification work converts six months of building into exam readiness: domains mapped to labs, two timed mocks, remediation by rebuilding rather than re-reading \cite{databricksCertification}.

The mental model: \textbf{you are ready when the platform redeploys without you}. The clean-clone drill---Terraform apply, bundle deploy, pipelines green, dashboards live, monitors armed---is the complete answer to ``do you understand this system?'' Everything else this week is measurement and polish.

\begin{definitionbox}
\textbf{Certification readiness} (this program's definition): two consecutive timed mock exams at 80\%+ with no domain below 70\%, plus a capstone that redeploys from a clean clone and survives a 15-minute architecture defense with cost, latency, and volume numbers from memory.
\end{definitionbox}

\begin{keyterms}
\textbf{Domain}: a scored exam section (Databricks platform, Delta Lake, ETL with Spark/DLT, production pipelines, governance). \textbf{Mock exam}: a timed, closed-book practice run scored by domain. \textbf{Remediation}: rebuilding the lab that exercises a weak domain. \textbf{Clean-clone drill}: full redeploy from the repository into a fresh workspace. \textbf{Architecture defense}: the narrated trade-off review of the final platform.
\end{keyterms}

\section{Monday: The Integrated Architecture Review}
\subsection{One Map}
Draw the whole platform on one page before touching any exam material: ingestion (Auto Loader, Kafka), medallion storage (Delta, liquid clustered), pipelines (DLT + Workflows), serving (SQL warehouse, AI/BI, Genie), ML (features, registry, endpoints, monitors), governance (UC, lineage, audit, sharing), and delivery (bundles, Terraform, Actions). For every box: one sentence of purpose, one number (volume, latency, or cost), one known limitation. That page is your capstone diagram, your exam warm-up, and your interview artifact.

\subsection{The Weak-Link Audit}
Walk the map asking ``what breaks this?'' per component: checkpoint deletion, skewed key, stale feature job, expired token, drifted schema, runaway DBU spend. Each answer should be a mechanism you built (alert, gate, runbook)---any answer of ``I'd notice'' is a gap to close before the defense.

\section{Tuesday: Certification Domains}
\subsection{Mapping Domains to Chapters}
The Databricks Data Engineer exams weight: the platform and tooling, Delta Lake, ETL with Spark (including DLT and Auto Loader), production pipelines and orchestration, and governance with Unity Catalog \cite{databricksCertification}. Map each domain to its chapters (1/4, 2/5, 9--11, 12--16, 8/21/22) and score your confidence per domain from mock evidence, not feeling.

\begin{table}[H]
\centering
\small
\begin{tabular}{@{}p{4.6cm}p{3.6cm}p{4.6cm}@{}}
\toprule
\textbf{Exam domain} & \textbf{Chapters} & \textbf{Rebuild these labs if weak} \\
\midrule
Databricks platform \& compute & 0, 1, 6 & Cluster policies, warehouse sizing lab \\
Delta Lake & 2, 5 & Enforcement/evolution/RESTORE drill; optimization benchmark \\
ETL with Spark \& DLT & 4, 9, 10 & Skew remediation; expectations poison drill \\
Incremental ingestion \& CDC & 11 & Auto Loader drift drill; SCD2 flow \\
Production pipelines & 12, 23 & Repair-run drill; bundle CI/CD lab \\
Streaming & 13--16 & Restart proof; watermark/dedup lab \\
Governance \& security & 8, 21, 22 & Masking + negative tests; share revocation \\
ML awareness & 17--20 & Registry alias flow; drift drill \\
\bottomrule
\end{tabular}
\caption{Exam domains mapped to chapters and their remedial labs.}
\label{tab:ch24-domains}
\end{table}

\section{Wednesday: Mock Exam Discipline}
\subsection{How to Take the Mocks}
Timed, closed-book, no notes, scored by domain. The diagnostic value is in the \emph{misses}: classify each as knowledge gap (never learned it), reasoning error (misread the scenario), or time management. Knowledge gaps remediate by rebuilding the lab in \cref{tab:ch24-domains}; reasoning errors by re-reading two scenario questions slowly; time issues by question budgeting. Reviewing answer keys without rebuilding the lab converts a mock into trivia.

\subsection{The Two-Pass Standard}
Readiness is two consecutive 80\%+ with no domain under 70\% \cite{databricksCertification}. One good mock is variance; two is signal. Book the exam when the standard is met---and keep the lab-rebuild habit; it is also how senior engineers stay honest.

\begin{pitfall}
\textbf{Studying features instead of failure modes.} Scenario questions reward knowing what breaks: which operation needs a new checkpoint, when expectations fail a pipeline, what \texttt{VACUUM} destroys, why a shared job cluster corrupts attribution. The labs taught these by breaking things---the exam tests exactly that.
\end{pitfall}

\section{Thursday: Final Capstone Build}
See \cref{sec:ch24-capstone} for the full specification. Today is the build day: the clean-clone drill, the replay drill, the evidence pack assembly.

\section{Friday Use Case: The Architecture Defense}
Present the platform in 15 minutes as if to a staff engineer who has never seen it.
\index{architecture defense}\index{presentation}

\subsection{The Structure That Scores}
Problem and scale first (``60 GB nightly plus 5k events/second, 200 analysts''), then the diagram, then two design decisions with trade-offs narrated (liquid clustering versus partitioning; triggered versus continuous), one failure story (what broke, how detected, what changed), and the closing numbers (DBU cost per month, p95 latency, rows per day). The examiner's follow-ups will target your stated trade-offs---choose ones you can defend two levels deep.

\begin{pitfall}
\textbf{The feature tour.} Listing services is not architecture. Every box must earn its place with a requirement it serves; every requirement must land in a box. If a box has no requirement, remove it; if a requirement has no box, that is the gap the defense exists to find.
\end{pitfall}

\subsection{Risks and Mitigations}
\begin{table}[H]
\centering
\small
\begin{tabular}{@{}p{4.6cm}p{7.6cm}@{}}
\toprule
\textbf{Risk} & \textbf{Mitigation} \\
\midrule
Capstone works only with owner's credentials & Clean-clone drill runs as the CI service principal \\
Evidence pack rots after exam week & Drills and evidence refresh quarterly, owned by the team \\
Mock-exam cramming without retention & Readiness standard requires consecutive passes weeks apart \\
Architecture defense reveals an unknown box & One-page map review with a peer before the presentation \\
Portfolio setup steps drift from reality & A stranger follows the README cold before publishing \\
\bottomrule
\end{tabular}
\caption{Week 24 scenario risks and mitigations.}
\label{tab:ch24-risks}
\end{table}

\section{Design Decision Guide}
\index{decision guide!Week 24}%
Certification and capstone decisions are about evidence: what you can show beats what you can say.

\begin{table}[H]
\centering
\small
\begin{tabular}{@{}p{3.4cm}p{4.4cm}p{4.6cm}@{}}
\toprule
\textbf{Decision} & \textbf{Choose the first when} & \textbf{Choose the second when} \\
\midrule
Lab rebuild vs.\ re-reading notes & Any domain below 70\% on a mock & Reviewing already-strong domains \\
Deep defense vs.\ feature tour & The 15-minute review & Never the tour \\
Associate vs.\ Professional exam & First certification, breadth & Proven production experience, depth \\
Scope cut vs.\ schedule slip & Capstone weeks & --- (protect evidence quality) \\
Book exam at standard vs.\ wait & Two consecutive 80\%+ mocks & One good mock---variance, not signal \\
Public portfolio vs.\ private repos & Job hunting & Employer constraints \\
\bottomrule
\end{tabular}
\caption{Week 24 decision guide.}
\label{tab:ch24-decisions}
\end{table}

The week's real deliverable is a habit: evidence over assertion, drills over documents, rebuilds over re-reads. The certification is the receipt; the platform is the proof.

\section{Operational Guidance for Week 24}
\index{operational guidance!Week 24}%
\subsection{Performance and Reliability}
The evidence pack is a living system: mock scores logged per domain, clean-clone drill timings trended per run, and the runbook re-executed (not re-read) quarterly. A clean-clone drill that takes longer each quarter is accumulating hand-rolled state---treat drill-duration regression like test-duration regression. Keep the defense recording current; the platform changes and the 15-minute story must change with it.

\subsection{Security and Governance Checklist}
\begin{itemize}
    \item Evidence packs (lineage, audit, cost) stored in the governed catalog, access-controlled.
    \item Clean-clone drill runs in an isolated workspace with no production credentials.
    \item Mock-exam artifacts contain no real platform data or secrets.
    \item The final platform's rollback and DR runbooks are attached to the release.
\end{itemize}

\subsection{Troubleshooting Playbook}
\begin{table}[H]
\centering
\small
\begin{tabular}{@{}p{3.6cm}p{4.2cm}p{4.8cm}@{}}
\toprule
\textbf{Symptom} & \textbf{Likely cause} & \textbf{First action} \\
\midrule
Mock score plateaus & Reviewing answers instead of rebuilding labs & Rebuild the weakest domain's lab from scratch \\
Clean-clone drill fails & Hand-rolled state not in code & Add the missing resource to bundle/Terraform; rerun \\
Defense runs over time & Feature tour instead of decisions & Cut to two trade-offs; rehearse the 90-second version \\
Drill passes but slowly & Cold-start and quota friction & Pre-warm quotas; document the critical path timings \\
\bottomrule
\end{tabular}
\caption{Week 24 troubleshooting quick reference.}
\label{tab:ch24-trouble}
\end{table}

\section{Interview Q\&A}
\subsection{Capstone-Grade Questions}
\begin{enumerate}
    \item \textbf{Q: Walk me through your platform end to end.}\\
    A: Practice the 90-second version: sources $\rightarrow$ ingestion $\rightarrow$ medallion $\rightarrow$ serving/ML $\rightarrow$ governance $\rightarrow$ delivery, with one number per arrow. If it takes five minutes, you do not yet own it.
    \item \textbf{Q: What would you rebuild differently at 10x scale?}\\
    A: Have the real answer: usually partition/parallelism ceilings, the orchestration fan-in, and the cost model. ``Nothing'' fails the question.
    \item \textbf{Q: Your most expensive mistake in this build?}\\
    A: The untested rollback, the unmonitored endpoint, the leaked token---name one, quantify it, show the mechanism that now prevents it.
    \item \textbf{Q: How does your platform prove data correctness to a skeptic?}\\
    A: Reconciliation arithmetic per transition, expectation metrics from the event log, replay drills, and the audit trail---evidence, not assurances.
    \item \textbf{Q: Why Databricks for this, over a point-solution stack?}\\
    A: One governed storage and catalog under batch, streaming, SQL, and ML---the alternative is four governance models and four copies. Know the honest counter-costs (DBU economics, Spark operational depth) too.
\end{enumerate}

\section{Further Reading}
The certification page lists current exams, domains, and preparation paths \cite{databricksCertification}; Databricks Academy self-paced courses align with the domains. Revisit the weakest domain's chapter Further Reading: platform \cite{databricksDocs}, Delta \cite{armbrust2020delta,deltaLakeDocs}, Spark \cite{sparkDocs,chambersZaharia}, pipelines \cite{databricksDLT,databricksWorkflows}, governance \cite{databricksUnityCatalog}, ML \cite{databricksMLflow,databricksMosaicAI}.

\section*{Key Takeaways}
\begin{itemize}
    \item Readiness = two consecutive 80\%+ mocks with no domain under 70\% + a clean-clone redeploy.
    \item Remediate by rebuilding labs, not re-reading notes.
    \item The one-page map with a purpose, a number, and a limitation per box is the defense and the study guide.
    \item The defense scores trade-off reasoning: two decisions defended deeply beat ten features listed.
    \item The portfolio README---problem, diagram, setup, evidence, cost, lessons---is the interview artifact.
\end{itemize}

\section{Exercises}
\begin{enumerate}
    \item \textbf{(Conceptual)} Write your 90-second platform walkthrough; time it; cut it to 90 seconds.
    \item \textbf{(Conceptual)} For each \cref{tab:ch24-domains} domain, name its three most failure-mode-flavored facts.
    \item \textbf{(Hands-on)} Run both mocks under exam conditions; log scores by domain.
    \item \textbf{(Hands-on)} Rebuild the lab for your weakest domain from scratch, in a fresh schema, without your notes.
    \item \textbf{(Hands-on)} Execute the clean-clone drill in a fresh workspace; record timings per step.
    \item \textbf{(Design)} Write the ``10x scale'' change list for your capstone with one justification each.
    \item \textbf{(Operations)} Record the 15-minute defense; review it against the structure in Friday's section.
    \item \textbf{(Portfolio)} Publish all six capstone repos with the standard README structure; have a peer follow one setup cold.
\end{enumerate}

\section{Milestone 6 Capstone: The Governed Data Platform}\label{sec:ch24-capstone}
\index{Milestone 6 capstone}\index{capstone project!Part VI}\index{clean-clone drill}

\begin{capstonebox}
\textbf{Mission.} Deliver the complete program as one platform: medallion lakehouse with optimized Delta storage, DLT pipelines with expectations, Auto Loader + Kafka streaming, an MLflow-managed model served and monitored, Unity Catalog governance with lineage/audit/sharing---all defined as code, promoted dev$\rightarrow$staging$\rightarrow$prod by CI, with cost and security evidence attached.

\textbf{Estimated time.} 12--16 hours across the week.

\textbf{What you'll practice.}
\begin{itemize}
    \item Extracting five months of hand-built artifacts into bundle/Terraform code.
    \item The clean-clone redeploy drill as the definitive correctness proof.
    \item Governance, cost, and reliability evidence packs.
    \item The certification-style architecture defense.
\end{itemize}
\end{capstonebox}

\subsection*{Scenario}
Your capstone platform from Milestones 1--5 works---in your workspace, with your credentials, by your hands. Meridian Retail's platform team now requires it to be \emph{operable by strangers}: deployable from the repo, governed by default, costed per domain, and defensible in an architecture review. This milestone closes that gap.

\subsection*{Requirements}
\begin{itemize}
    \item One bundle (jobs, pipelines, dashboards, serving endpoint, monitors) with dev/staging/prod targets; Terraform for the perimeter.
    \item Clean-environment redeploy: \texttt{terraform apply} + \texttt{databricks bundle deploy -t prod} rebuilds the platform from the repo alone.
    \item Batch and streaming paths reconciled in Gold; the replay drill producing zero duplicates.
    \item Governance evidence: access matrix implemented, lineage report, PII audit alert, one Delta Sharing share with revocation test.
    \item CI/CD with PR validation, staging gate, approved prod promotion, and a documented rollback.
    \item Cost report per domain with two implemented optimizations; security review checklist complete.
    \item Two mock exams meeting the readiness standard; the recorded 15-minute defense.
\end{itemize}

\subsection*{Reference Architecture}
\begin{figure}[H]
    \centering
    \resizebox{\textwidth}{!}{%
    \begin{tikzpicture}[
        box/.style={rectangle, rounded corners, draw=primary, thick, fill=primary!7, align=center, minimum width=2.6cm, minimum height=0.95cm, font=\small\bfseries},
        gov/.style={rectangle, rounded corners, draw=red!60!black, thick, fill=red!5, align=center, minimum width=3.0cm, minimum height=0.85cm, font=\scriptsize\bfseries},
        ops/.style={rectangle, rounded corners, draw=teal!55!black, thick, fill=teal!8, align=center, minimum width=3.0cm, minimum height=0.85cm, font=\scriptsize\bfseries},
        flow/.style={-{Latex[length=2mm]}, draw=primary, thick}
    ]
        \node[box] (src) at (0,0) {Files + Kafka\\+ CDC};
        \node[box] (lake) at (3.4,0) {Medallion\\(DLT +\\expectations)};
        \node[box] (serve) at (6.8,0) {Serving\\SQL + AI/BI\\+ Genie};
        \node[box] (ml) at (10.2,0) {ML\\features, registry,\\endpoint, monitors};
        \node[gov] (uc) at (3.4,-1.7) {Unity Catalog\\grants, lineage,\\audit, sharing};
        \node[ops] (cicd) at (8.5,-1.7) {Bundles + Terraform\\GitHub Actions\\dev$\rightarrow$stg$\rightarrow$prod};
        \node[ops] (fin) at (12.6,-1.7) {FinOps\\chargeback};
        \draw[flow] (src) -- (lake);
        \draw[flow] (lake) -- (serve);
        \draw[flow] (lake) -- (ml);
        \draw[flow, dashed] (uc) -- (lake);
        \draw[flow, dashed] (cicd) -- (serve);
        \draw[flow, dashed] (fin) -- (ml);
    \end{tikzpicture}%
    }
    \caption{Milestone 6: the complete platform---one governed estate, deployed as code, measured continuously.}
    \label{fig:ch24-capstone-arch}
\end{figure}

\subsection*{Implementation Steps}
\begin{enumerate}
    \item Inventory every hand-built artifact from Milestones 1--5; express each as bundle YAML or Terraform.
    \item Define targets with variables; wire the Actions pipeline with the staging gate.
    \item Run the clean-clone drill in a fresh workspace; fix everything the drill finds.
    \item Assemble the evidence packs: governance, reliability (drills), cost, security.
    \item Both mocks, the remediation rebuilds, and the recorded defense.
\end{enumerate}

\subsection*{Deliverables}
\begin{itemize}
    \item The repository: bundles, Terraform, pipelines, tests, README per the portfolio guide.
    \item The drill logs: clean-clone timings, replay-drill counts, revocation evidence.
    \item The cost and security packs; the mock scores; the defense recording.
\end{itemize}

\subsection*{Acceptance Criteria}
\begin{itemize}
    \item A teammate (or you, in a fresh workspace) redeploys the platform from the README alone.
    \item No human holds prod deploy rights; the pipeline audit trail is complete.
    \item Every acceptance criterion from Milestones 1--5 still passes against the code-deployed estate.
    \item The readiness standard (two mocks, 80\%+, no domain $<$70\%) is met.
\end{itemize}

\subsection*{Stretch Goals}
\begin{itemize}
    \item Add the drift-detection job proving prod matches the bundle registry.
    \item Publish one Gold dataset via Delta Sharing to a second (free) workspace with the revocation drill.
    \item Write the one-page ``operating this platform'' runbook a stranger could follow oncall.
\end{itemize}

\vspace{4mm}
\begin{center}
    \textit{``Consistency is the architecture of mastery.''}\\[1mm]
    You have built the platform. Now go defend it.
\end{center}
